% GENERATED FILE. Edit canonical lessons, then run npm run generate:books.
% canonical-source-hash: fnv1a64:910dad1b1884abe5
% canonical-lessons: ML-C569-campyan, ML-C569-etirali, ML-C569-turnamenru, ML-C569-munneruka, ML-C569-pinnilakuka

\chapter{Winning and Losing}
\label{ch:ml-winning-and-losing}
\hlchaptermodality{569}

\begin{chapteropening}
\textbf{By the end of this chapter:} \emph{I can say a sports champion, an opponent, a tournament, to advance and to fall behind.}

Complete the last lesson of chapter 569: I can say a sports champion, an opponent, a tournament, to advance and to fall behind.
\end{chapteropening}

\section[\texorpdfstring{\ml{ചാമ്പ്യൻ} (cāmpyan)}{ചാമ്പ്യൻ (cāmpyan)}]{\ml{ചാമ്പ്യൻ} (cāmpyan) — a sports champion}

\label{lesson:ML-C569-campyan}



\begin{quote}
Before the new one: say the Malayalam for a referee, then the Malayalam for to compete.
\end{quote}

\subsection*{You'll want to know: \ml{ചാമ്പ്യൻ}}
\textbf{\ml{ചാമ്പ്യൻ}} — \emph{cāmpyan} — \textquotedblleft{}a sports champion\textquotedblright{}.

\begin{quote}
\ml{ഞങ്ങളുടെ} \ml{സ്കൂൾ} \ml{ടീം} \ml{ജില്ലയിലെ} \ml{ചാമ്പ്യൻ} \ml{ആയി}.
\end{quote}

\begin{quote}
\emph{ñaṅṅaḷuṭe skūḷ ṭīṁ jillayile cāmpyan āyi.}
\end{quote}

\begin{quote}
Our school team became the district champion.
\end{quote}

Plural: \ml{ചാമ്പ്യന്മാർ}.

\begin{cousinweb}
\textbf{\ml{ചാമ്പ്യൻ}} (\emph{cāmpyan}) was borrowed from English \emph{champion}.
\end{cousinweb}

\begin{grammarlens}[title={What you've built}]
One more word for describing what you lived and felt.
\end{grammarlens}

\subsection*{Your turn}
\begin{itemize}
\raggedright
  \item \emph{Say it:} \emph{cāmpyan}
  \item \emph{Say it:} \emph{cāmpyan}, once more
  \item \emph{Say it:} \emph{ñaṅṅaḷuṭe skūḷ ṭīṁ jillayile cāmpyan āyi.} — Our school team became the district champion.
  \item \emph{Recall:} say the Malayalam for a referee, then the Malayalam for to compete, then say \emph{cāmpyan} again
\end{itemize}

\begin{tcolorbox}[breakable,colback=teal!4,colframe=teal!35!black,title={Before you move on}]
What does \ml{ചാമ്പ്യൻ} mean? (\textquotedblleft{}A sports champion\textquotedblright{}.) Say it once more.
\end{tcolorbox}
\section[\texorpdfstring{\ml{എതിരാളി} (etirāḷi)}{എതിരാളി (etirāḷi)}]{\ml{എതിരാളി} (etirāḷi) — an opponent; a rival}

\label{lesson:ML-C569-etirali}



\begin{quote}
Before the new one: say the Malayalam for to compete, then the Malayalam for a sports champion.
\end{quote}

\subsection*{You'll want to know: \ml{എതിരാളി}}
\textbf{\ml{എതിരാളി}} — \emph{etirāḷi} — \textquotedblleft{}an opponent; a rival\textquotedblright{}.

\begin{quote}
\ml{അവസാന} \ml{കളിയിൽ} \ml{ഞങ്ങളുടെ} \ml{എതിരാളി} \ml{ശക്തമായ} \ml{ടീമായിരുന്നു}.
\end{quote}

\begin{quote}
\emph{avasāna kaḷiyil ñaṅṅaḷuṭe etirāḷi śaktamāya ṭīmāyirunnu.}
\end{quote}

\begin{quote}
In the last match our opponents were a strong team.
\end{quote}

\ml{എതിർ} opposite + \ml{ആൾ} a person; \ml{എതിരെ}, against, is already familiar.

\begin{cousinweb}
\textbf{\ml{എതിരാളി}} (\emph{etirāḷi}) is native Dravidian, a compound of \ml{എതിർ}, \textquotedblleft{}opposite, adverse\textquotedblright{} (DEDR 795), and \ml{ആൾ}, \textquotedblleft{}a person\textquotedblright{} (DEDR 399).
\end{cousinweb}

\begin{grammarlens}[title={What you've built}]
One more word for describing what you lived and felt.
\end{grammarlens}

\subsection*{Your turn}
\begin{itemize}
\raggedright
  \item \emph{Say it:} \emph{etirāḷi}
  \item \emph{Say it:} \emph{etirāḷi}, once more
  \item \emph{Say it:} \emph{avasāna kaḷiyil ñaṅṅaḷuṭe etirāḷi śaktamāya ṭīmāyirunnu.} — In the last match our opponents were a strong team.
  \item \emph{Recall:} say the Malayalam for to compete, then the Malayalam for a sports champion, then say \emph{etirāḷi} again
\end{itemize}

\begin{tcolorbox}[breakable,colback=teal!4,colframe=teal!35!black,title={Before you move on}]
What does \ml{എതിരാളി} mean? (\textquotedblleft{}An opponent; a rival\textquotedblright{}.) Say it once more.
\end{tcolorbox}
\section[\texorpdfstring{\ml{ടൂർണമെന്റ്} (ṭūrṇamenṟŭ)}{ടൂർണമെന്റ് (ṭūrṇamenṟŭ)}]{\ml{ടൂർണമെന്റ്} (ṭūrṇamenṟŭ) — a tournament}

\label{lesson:ML-C569-turnamenru}



\begin{quote}
Before the new one: say the Malayalam for a sports champion, then the Malayalam for an opponent.
\end{quote}

\subsection*{You'll want to know: \ml{ടൂർണമെന്റ്}}
\textbf{\ml{ടൂർണമെന്റ്}} — \emph{ṭūrṇamenṟŭ} — \textquotedblleft{}a tournament\textquotedblright{}.

\begin{quote}
\ml{അടുത്ത} \ml{മാസം} \ml{സ്കൂളിൽ} \ml{ഒരു} \ml{ഫുട്ബോൾ} \ml{ടൂർണമെന്റ്} \ml{ഉണ്ട്}.
\end{quote}

\begin{quote}
\emph{aṭutta māsaṁ skūḷil oru phuṭbōḷ ṭūrṇamenṟŭ uṇṭŭ.}
\end{quote}

\begin{quote}
There is a football tournament at school next month.
\end{quote}

English -nt comes out as -\ml{ന്റ്}, as in \ml{പ്രിന്റർ}, a printer.

\begin{cousinweb}
\textbf{\ml{ടൂർണമെന്റ്}} (\emph{ṭūrṇamenṟŭ}) was borrowed from English \emph{tournament}.
\end{cousinweb}

\begin{grammarlens}[title={What you've built}]
One more word for describing what you lived and felt.
\end{grammarlens}

\subsection*{Your turn}
\begin{itemize}
\raggedright
  \item \emph{Say it:} \emph{ṭūrṇamenṟŭ}
  \item \emph{Say it:} \emph{ṭūrṇamenṟŭ}, once more
  \item \emph{Say it:} \emph{aṭutta māsaṁ skūḷil oru phuṭbōḷ ṭūrṇamenṟŭ uṇṭŭ.} — There is a football tournament at school next month.
  \item \emph{Recall:} say the Malayalam for a sports champion, then the Malayalam for an opponent, then say \emph{ṭūrṇamenṟŭ} again
\end{itemize}

\begin{tcolorbox}[breakable,colback=teal!4,colframe=teal!35!black,title={Before you move on}]
What does \ml{ടൂർണമെന്റ്} mean? (\textquotedblleft{}A tournament\textquotedblright{}.) Say it once more.
\end{tcolorbox}
\section[\texorpdfstring{\ml{മുന്നേറുക} (munnēṟuka)}{മുന്നേറുക (munnēṟuka)}]{\ml{മുന്നേറുക} (munnēṟuka) — to advance; to forge ahead}

\label{lesson:ML-C569-munneruka}



\begin{quote}
Before the new one: say the Malayalam for an opponent, then the Malayalam for a tournament.
\end{quote}

\subsection*{You'll want to know: \ml{മുന്നേറുക}}
\textbf{\ml{മുന്നേറുക}} — \emph{munnēṟuka} — \textquotedblleft{}to advance; to forge ahead\textquotedblright{}.

\begin{quote}
\ml{രണ്ടാം} \ml{പകുതിയിൽ} \ml{ഞങ്ങളുടെ} \ml{ടീം} \ml{മുന്നേറി}.
\end{quote}

\begin{quote}
\emph{raṇṭāṁ pakutiyil ñaṅṅaḷuṭe ṭīṁ munnēṟi.}
\end{quote}

\begin{quote}
In the second half our team pulled ahead.
\end{quote}

Past: \ml{മുന്നേറി}, \emph{munnēṟi}.

\ml{മുൻ} front + \ml{ഏറുക} to rise, to climb: to move to the front.

\begin{cousinweb}
\textbf{\ml{മുന്നേറുക}} (\emph{munnēṟuka}) is native Dravidian, a compound of \ml{മുൻ}, \textquotedblleft{}front, before\textquotedblright{} (DEDR 5020), and \ml{ഏറുക}, \textquotedblleft{}to rise, to ascend\textquotedblright{} (DEDR 916).
\end{cousinweb}

\begin{grammarlens}[title={What you've built}]
One more word for describing what you lived and felt.
\end{grammarlens}

\subsection*{Your turn}
\begin{itemize}
\raggedright
  \item \emph{Say it:} \emph{munnēṟuka}
  \item \emph{Say it:} \emph{munnēṟuka}, once more
  \item \emph{Say it:} \emph{raṇṭāṁ pakutiyil ñaṅṅaḷuṭe ṭīṁ munnēṟi.} — In the second half our team pulled ahead.
  \item \emph{Recall:} say the Malayalam for an opponent, then the Malayalam for a tournament, then say \emph{munnēṟuka} again
\end{itemize}

\begin{tcolorbox}[breakable,colback=teal!4,colframe=teal!35!black,title={Before you move on}]
What does \ml{മുന്നേറുക} mean? (\textquotedblleft{}To advance; to forge ahead\textquotedblright{}.) Say it once more.
\end{tcolorbox}
\section[\texorpdfstring{\ml{പിന്നിലാകുക} (pinnilākuka)}{പിന്നിലാകുക (pinnilākuka)}]{\ml{പിന്നിലാകുക} (pinnilākuka) — to fall behind}

\label{lesson:ML-C569-pinnilakuka}



\begin{quote}
Before the new one: say the Malayalam for a tournament, then the Malayalam for to advance.
\end{quote}

\subsection*{You'll want to know: \ml{പിന്നിലാകുക}}
\textbf{\ml{പിന്നിലാകുക}} — \emph{pinnilākuka} — \textquotedblleft{}to fall behind\textquotedblright{}.

\begin{quote}
\ml{ആദ്യ} \ml{പകുതിയിൽ} \ml{ഞങ്ങളുടെ} \ml{ടീം} \ml{രണ്ട്} \ml{ഗോളിന്} \ml{പിന്നിലായി}.
\end{quote}

\begin{quote}
\emph{ādya pakutiyil ñaṅṅaḷuṭe ṭīṁ raṇṭŭ gōḷinŭ pinnilāyi.}
\end{quote}

\begin{quote}
In the first half our team fell two goals behind.
\end{quote}

Past: \ml{പിന്നിലായി}, \emph{pinnilāyi}.

\ml{പിന്നിൽ} behind + \ml{ആകുക} to become.

\begin{cousinweb}
\textbf{\ml{പിന്നിലാകുക}} (\emph{pinnilākuka}) is native Dravidian, built from \ml{പിന്നിൽ}, \textquotedblleft{}behind\textquotedblright{}, from \ml{പിൻ} (DEDR 4205), and \ml{ആകുക}, \textquotedblleft{}to become\textquotedblright{}.
\end{cousinweb}

\begin{grammarlens}[title={What you've built}]
One more word for describing what you lived and felt.
\end{grammarlens}

\subsection*{Your turn}
\begin{itemize}
\raggedright
  \item \emph{Say it:} \emph{pinnilākuka}
  \item \emph{Say it:} \emph{pinnilākuka}, once more
  \item \emph{Say it:} \emph{ādya pakutiyil ñaṅṅaḷuṭe ṭīṁ raṇṭŭ gōḷinŭ pinnilāyi.} — In the first half our team fell two goals behind.
  \item \emph{Recall:} say the Malayalam for a tournament, then the Malayalam for to advance, then say \emph{pinnilākuka} again
\end{itemize}

\begin{tcolorbox}[breakable,colback=teal!4,colframe=teal!35!black,title={Before you move on}]
What does \ml{പിന്നിലാകുക} mean? (\textquotedblleft{}To fall behind\textquotedblright{}.) Say it once more.
\end{tcolorbox}
